\documentclass[11pt]{article}
\usepackage[utf8]{inputenc}
\usepackage{amsmath,amssymb}
\usepackage{hyperref}
\title{Efficient Organic Photovoltaic Materials: A Resource-Aware Approach}
\author{Anonymous}
\date{}
\begin{document}
\maketitle

\begin{abstract}
This paper studies Organic Photovoltaic Materials. Grounded in a real, citable literature \cite{ref1,ref2,ref3,ref4,ref5,ref6,ref7,ref8},
we propose a focused method and report \textbf{illustrative} experiments.
All references below are real and verifiable (OpenAlex/arXiv); the reported
numbers are simulated to illustrate intended behaviour.
\end{abstract}

\section{Introduction}
Organic Photovoltaic Materials is an active research area. This work targets a concrete gap identified from
the recent literature and contributes a method together with an evaluation protocol.

\section{Related Work (real literature)}
The following works, retrieved from public bibliographic APIs, frame our study:
\begin{itemize}
\item Burschka et al. (2013) — "Sequential deposition as a route to high-performance perovskite-sensitized solar cells", Nature (cited 9474).
\item Skotheim et al. (1986) — "Handbook of conducting polymers", M. Dekker eBooks (cited 8071).
\item Green et al. (2014) — "The emergence of perovskite solar cells", Nature Photonics (cited 7340).
\item Xing et al. (2013) — "Long-Range Balanced Electron- and Hole-Transport Lengths in Organic-Inorganic CH <sub>3</sub> NH <sub>3</sub> PbI <sub>3</sub>", Science (cited 6786).
\item Jeon et al. (2014) — "Solvent engineering for high-performance inorganic–organic hybrid perovskite solar cells", Nature Materials (cited 6667).
\item Jeon et al. (2015) — "Compositional engineering of perovskite materials for high-performance solar cells", Nature (cited 6398).
\end{itemize}

\section{Method}
We reduce the dominant computational cost via adaptive resource allocation, activating only the components needed per input. We instantiate this idea for Organic Photovoltaic Materials and analyse its cost/quality trade-off.

\section{Experiments (Illustrative)}
\begin{center}
\begin{tabular}{lcc}
System & Primary Metric & Efficiency \\ \hline
Strong baseline & 0.812 & 1.00$\times$ \\
Ablation & 0.828 & 0.92$\times$ \\
\textbf{Ours} & \textbf{0.851} & \textbf{0.71$\times$}
\end{tabular}
\end{center}
\emph{Metrics simulated to illustrate the intended effect; not measured.}

\section{Conclusion}
We connected a real literature to a concrete method for Organic Photovoltaic Materials. Future work will
validate the illustrative results with real experiments.

\begin{thebibliography}{9}
\bibitem{ref1} Julian Burschka, Norman Pellet, Soo‐Jin Moon et al.. Sequential deposition as a route to high-performance perovskite-sensitized solar cells. \emph{Nature}, 2013. DOI:10.1038/nature12340
\bibitem{ref2} Terje A. Skotheim, Ronald L. Elsenbaumer, John R. Reynolds. Handbook of conducting polymers. \emph{M. Dekker eBooks}, 1986. 
\bibitem{ref3} Martin A. Green, Anita Ho‐Baillie, Henry J. Snaith. The emergence of perovskite solar cells. \emph{Nature Photonics}, 2014. DOI:10.1038/nphoton.2014.134
\bibitem{ref4} Guichuan Xing, Nripan Mathews, Shuangyong Sun et al.. Long-Range Balanced Electron- and Hole-Transport Lengths in Organic-Inorganic CH <sub>3</sub> NH <sub>3</sub> PbI <sub>3</sub>. \emph{Science}, 2013. DOI:10.1126/science.1243167
\bibitem{ref5} Nam Joong Jeon, Jun Hong Noh, Young Chan Kim et al.. Solvent engineering for high-performance inorganic–organic hybrid perovskite solar cells. \emph{Nature Materials}, 2014. DOI:10.1038/nmat4014
\bibitem{ref6} Nam Joong Jeon, Jun Hong Noh, Woon Seok Yang et al.. Compositional engineering of perovskite materials for high-performance solar cells. \emph{Nature}, 2015. DOI:10.1038/nature14133
\bibitem{ref7} C. W. Tang. Two-layer organic photovoltaic cell. \emph{Applied Physics Letters}, 1986. DOI:10.1063/1.96937
\bibitem{ref8} . The Physics of Semiconductor Devices. \emph{Springer proceedings in physics}, 2024. DOI:10.1007/978-981-97-1571-8
\end{thebibliography}
\end{document}
